\documentclass[../main]{subfiles}
\ifSubfilesClassLoaded{\setcounter{chapter}{1}}{}
\begin{document}

\chapter{Landasan Teori dan Konsep Dasar Logika Matematika}

Setelah menyelesaikan bab ini, mahasiswa diharapkan tidak lagi memandang logika matematika hanya sebagai hafalan teoretis. Bab ini memaparkan perkembangan logika dari fondasi filosofis hingga penerapannya dalam pengembangan perangkat lunak, kecerdasan buatan, dan rangkaian terpadu.

\begin{subcpmk}
  \item \textbf{Sub-CPMK 1.1:} Mahasiswa mampu menjelaskan konsep dasar proposisi, nilai kebenaran, dan kata hubung logika (konjungsi, disjungsi, negasi, implikasi, biimplikasi).
  \item \textbf{Sub-CPMK 1.2:} Mahasiswa mampu menggunakan tabel kebenaran untuk mengevaluasi pernyataan majemuk serta mengidentifikasi tautologi, kontradiksi, dan kontingensi.
  \item \textbf{Sub-CPMK 1.3:} Mahasiswa mampu menganalisis validitas argumen menggunakan aturan inferensi (modus ponens, modus tollens, silogisme) dalam memecahkan persoalan komputasi dasar.
  \item \textbf{Sub-CPMK 2.1:} Mahasiswa mampu menjelaskan konsep dasar logika predikat, fungsi proposisional, dan ruang lingkup pembicaraan (\textit{universe of discourse}).
  \item \textbf{Sub-CPMK 2.2:} Mahasiswa mampu menerjemahkan pernyataan bahasa alami ke dalam ekspresi logika predikat menggunakan kuantor universal dan eksistensial.
  \item \textbf{Sub-CPMK 2.3:} Mahasiswa mampu memodelkan masalah komputasi yang terstruktur dengan merepresentasikan logika predikat sebagai basis perancangan algoritma.
  \item \textbf{Sub-CPMK 3.1:} Mahasiswa mampu mendefinisikan prinsip-prinsip dan hukum-hukum dasar Aljabar Boolean.
  \item \textbf{Sub-CPMK 3.2:} Mahasiswa mampu menyederhanakan fungsi logika menggunakan hukum Aljabar Boolean atau Peta Karnaugh (K-Map).
  \item \textbf{Sub-CPMK 3.3:} Mahasiswa mampu mengonversi dan merepresentasikan fungsi logika yang telah disederhanakan ke dalam bentuk rangkaian gerbang logika dasar pembentuk sistem komputer.
  \item \textbf{Sub-CPMK 4.1:} Mahasiswa mampu menerapkan berbagai metode pembuktian matematis dasar (pembuktian langsung, pembuktian tidak langsung/kontradiksi, dan kontraposisi).
  \item \textbf{Sub-CPMK 4.2:} Mahasiswa mampu mengaplikasikan prinsip Induksi Matematika untuk membuktikan kebenaran sebuah deret atau properti dari algoritma rekursif sederhana.
  \item \textbf{Sub-CPMK 4.3:} Mahasiswa mampu menyusun argumen matematis formal untuk memvalidasi langkah-langkah dalam penyelesaian masalah komputasi.
\end{subcpmk}

% ============================================================
% MATERI POKOK
% ============================================================
\section{Akar Filosofis: Dari Silogisme hingga Logika Simbolik}

Sebelum mesin-mesin modern mampu memproses data dalam skala sangat besar, landasan ilmu komputer (*Computer Science*) telah dirumuskan terlebih dahulu oleh para filsuf kuno melalui gagasan \textbf{"Hukum Pasti Berpikir"} \cite{wiki_proplogic}.

\logic{Logika Matematika} merupakan instrumen untuk mengalihkan pernyataan dalam bahasa alamiah yang mengandung ambiguitas menjadi rumus yang deterministik, terdefinisi dengan jelas, dan dapat diproses oleh mesin \cite{rosen2019,copi2014,forallx}.

\subsection{Fajar Logika Yunani: Era Silogisme Aristoteles}

Pada rentang tahun 384--322 SM, Aristoteles merintis karya intelektual yang membakukan struktur perdebatan manusia menjadi rumus deduksi yang ketat, yang kemudian disebut \textbf{Silogisme} \cite{iep_logic_history}. Era ini menandai terobosan ketika keabsahan suatu pernyataan dinilai berdasarkan struktur penalarannya, terlepas dari kebenaran faktual isi pernyataan tersebut.

\begin{contoh}[Bentuk Klasik Silogisme Aristoteles]
\indent \textbf{Premis Mayor:} Semua mamalia wajib mengembuskan karbon dioksida. \\
\indent \textbf{Premis Minor:} Seekor kucing perak tergolong bangsa mamalia. \\
\indent \textbf{Validasi Kesimpulan:} Dengan demikian, kucing perak tersebut mengembuskan karbon dioksida.
\end{contoh}

Sumbangan bangsa Yunani Kuno (khususnya penerus filosofi Stoik) melahirkan dasar-dasar hubungan kondisional "Jika-Maka" (\textit{If-Then}). Namun, selama berabad-abad, analisis logika mereka belum dituangkan secara sistematis dalam bentuk matematis yang transparan.

\subsection{Era Revolusi Simbolik: George Boole \& Gottlob Frege}

Pada abad ke-19, logika berkembang menjadi disiplin yang berorientasi pada ketegasan penalaran matematis \cite{openlogic,copi2014}. Tokoh utamanya ialah matematikawan Inggris, \textbf{George Boole (1815-1864)}.

Dalam karyanya \textit{The Laws of Thought (1854)}, Boole mengemukakan gagasan bahwa nilai kebenaran suatu kalimat dapat direduksi menjadi rumus matematika, analog dengan operasi penjumlahan (+) dan perkalian ($\times$). Boole memperkenalkan \textbf{Aljabar Boolean}. Melalui gagasan ini, penalaran bahasa dapat dipandang sebagai operasi biner 1 (True) dan 0 (False).

Gagasan tersebut kemudian dikembangkan oleh \textbf{Gottlob Frege (1879)} melalui karyanya \textit{Begriffsschrift} \cite{copi2014}. Frege mengembangkan logika Boole dengan memperkenalkan \textbf{Logika Predikat}. Ia memperkenalkan kuantifier universal ``Semua'' ($\forall$) dan eksistensial ``Sebagian'' ($\exists$). Dengan demikian, pencarian kondisi majemuk pada basis data relasional (misalnya SQL) dapat dieksekusi secara presisi dalam sistem digital modern.

\begin{figure}[H]
\centering
\includegraphics[width=0.9\textwidth]{timeline-logika}
\caption{Timeline evolusi Logika Matematika: dari perkembangan pemikiran filsuf hingga penerapannya pada perangkat elektronik \cite{wiki_proplogic}.}
\label{fig:timeline-logika}
\end{figure}

\section{Anatomi Argumen: Premis, Konklusi, dan Penilaian Validitas}

Inti disiplin Logika Matematika bermuara pada kesanggupan untuk menentukan status sahih suatu argumen \cite{forallx,copi2014}. Mahasiswa diharapkan memahami bahwa pemetaan informasi ke dalam tiga entitas dilakukan melalui proposisi, aturan inferensi, dan kesimpulan.

\subsection{Komponen Pembangun Argumen Formal}

\begin{table}[H]
\centering
\small
\begin{tabular}{|>{\raggedright\arraybackslash}p{2.8cm}|>{\raggedright\arraybackslash}p{9.4cm}|}
\hline
\textbf{Terminologi} & \textbf{Anatomi Fungsional di Semesta Logika} \\
\hline
\textbf{1. Proposisi} & Satuan terkecil kalimat asertif yang tidak dapat dibelah; bernilai kebenaran BENAR (T) atau SALAH (F) \cite{rosen2019}. \\
\hline
\textbf{2. Premis} & Kumpulan hipotesis atau \textit{bukti} yang diasumsikan benar di awal untuk mengecek keutuhan pondasi argumen \cite{copi2014}. \\
\hline
\textbf{3. Kesimpulan (\textit{output})} & Kesimpulan yang dihasilkan dari penerapan aturan inferensi terhadap premis. \\
\hline
\textbf{4. Aturan Inferensi} & Aturan yang memverifikasi apakah langkah dari premis menuju kesimpulan mematuhi hukum logika \cite{iep_natded}. \\
\hline
\end{tabular}
\caption{Tinjauan anatomi argumen formal}
\label{tab:konsep-logika}
\end{table}

Secara visual, proses pembuktian teorema matematis dapat digambarkan sebagai alur penyusunan (Gambar~\ref{fig:struktur-argumen}). Premis diteruskan ke aturan inferensi. Jika aturan tersebut valid, kesimpulan yang dihasilkan adalah benar.

\begin{figure}[H]
\centering
\begin{tikzpicture}[
  node distance=1.2cm,
  premis/.style={rectangle, draw=black, fill=blue!15, rounded corners, minimum width=2.5cm, minimum height=0.8cm, font=\bfseries},
  inferensi/.style={rectangle, draw=black, fill=green!15, rounded corners, text width=3.2cm, minimum height=1cm, align=center, font=\bfseries},
  kesimpulan/.style={rectangle, draw=black, fill=orange!15, rounded corners, text width=2.3cm, minimum height=0.8cm, align=center, font=\bfseries},
  arrow/.style={->, thick, >=stealth}
]
  \node[premis] (p1) {Premis Induk 1 (Diasumsi T)};
  \node[premis, below=0.7cm of p1] (p2) {Premis Sub 2 (Diasumsi T)};
  \node[inferensi, right=2.5cm of $(p1)!0.5!(p2)$] (inf) {Proses \newline ATURAN INFERENSI};
  \node[kesimpulan, right=2.5cm of inf] (k) {Konklusi Akhir};
  \draw[arrow] (p1.east) -- (inf.west);
  \draw[arrow] (p2.east) -- (inf.west);
  \draw[arrow] (inf.east) -- (k.west);
\end{tikzpicture}
\caption{Visualisasi alur penyusunan argumen: dari premis menuju kesimpulan yang valid.}
\label{fig:struktur-argumen}
\end{figure}

\subsection{Perbedaan antara Fakta Empiris dan Validitas Aljabar}

Bagian yang paling sulit dalam pembelajaran logika adalah konsep validitas. \textbf{Suatu argumen dapat diberi predikat ``valid'' berdasarkan bentuk aljabarnya tanpa perlu selaras dengan fakta empiris maupun hukum fisika/biologi di dunia nyata.} \cite{forallx,copi2014}

\begin{definisi}[Hukum Absolut Argumen Valid]
Argumen dinyatakan \logic{valid} apabila struktur sintaksnya menjamin: jika semua premis diasumsikan Benar (T), maka tidak mungkin kesimpulan bernilai Salah (F). Kesahihan validitas ditentukan semata dari ketepatan jalur inferensinya, bukan dari kebenaran material premis.
\end{definisi}

\begin{contoh}[Valid Secara Struktur meskipun Isinya Tidak Realistis]
\textbf{Premis 1:} Jika Budi makan durian, maka Budi melahirkan alien. \\
\textbf{Premis 2:} Budi sedang makan durian. \\
\textbf{Kesimpulan:} Budi melahirkan alien. \\
\textbf{Penilaian: Argumen valid secara struktur.} \\
Penjelasan: Meskipun kesimpulan tidak sesuai fakta empiris (manusia tidak dapat melahirkan alien), struktur inferensi \textit{Modus Ponens} bekerja sesuai ketentuan. Jika (A) mengakibatkan (B) dan (A) benar, maka (B) harus benar. Validitas logika tidak bergantung pada kebenaran faktual premis.
\end{contoh}

\begin{contoh}[Argumen Terdengar Realistis namun Invalid---Fallacy of Affirming the Consequent]
\textbf{Premis 1:} Jika badai halilintar turun, maka stasiun TV mati. \\
\textbf{Premis 2:} Stasiun TV mati. \\
\textbf{Kesimpulan:} Badai halilintar turun. \\
\textbf{Penilaian: Argumen invalid.} \\
Penjelasan: Meskipun premis terdengar realistis, susunan argumen ini mengalami \textit{Fallacy of Affirming the Consequent} (Penegasan Konsekuen). Stasiun TV dapat mati karena sebab lain (misalnya kerusakan peralatan); kesimpulan tidak dijamin oleh premis. Aturan inferensinya tidak sesuai \cite{iep_natded}.
\end{contoh}
Mengkaji kesalahan inferensi seperti pada contoh di atas merupakan kompetensi penting bagi auditor sistem informasi saat menyusun spesifikasi persyaratan alur aplikasi.

\section{Logika Boolean dalam Relung Arsitektur Perangkat Lunak}

Beranjak dari konteks filosofis, bagian ini membahas penerapan logika langsung di layar terminal komputasi \cite{levin_discrete,saylor_cs202}. Pengembangan perangkat lunak di abad ke-21 bergantung pada \textbf{Aljabar Logika Boolean}.

\subsection{Boolean sebagai Dasar \textit{Control Flow}}

Struktur \textit{Control Flow} dan \textit{Branching Decision} (keputusan bercabang) pada setiap program ditentukan oleh evaluasi ekspresi logika Boolean, yang menyatakan nilai benar (1/\textit{True}) atau nilai salah (0/\textit{False}).

Tanpa evaluasi kondisi logika, program hanya mengeksekusi pernyataan secara berurutan tanpa percabangan. Evaluasi syarat logika (\texttt{IF}, \texttt{ELSE IF}, repetisi \texttt{WHILE}) memungkinkan terjadinya percabangan alur eksekusi. Mesin dapat mengelola alur berdasarkan hasil evaluasi \texttt{True} dan \texttt{False} \cite{levin_discrete}.

\begin{figure}[H]
\centering
\begin{tikzpicture}[
  node distance=0.8cm,
  decision/.style={diamond, draw=black, fill=red!15, aspect=2, align=center, font=\bfseries\small, text width=2.2cm},
  process/.style={rectangle, draw=black, fill=blue!15, rounded corners, minimum height=0.8cm, text width=2.8cm, align=center},
  start/.style={ellipse, draw=black, fill=green!20, font=\bfseries},
  arrow/.style={->, thick, >=stealth}
]
  \node[start] (start) {Mesin aktif (mulai)};
  \node[decision, below=of start] (cond) {Verifikasi Kondisi Boolean?};
  \node[process, below left=1cm and 1.2cm of cond] (t) {Eksekusi Cabang True};
  \node[process, below right=1cm and 1.2cm of cond] (f) {Lewati Cabang False};
  \node[process, below=2.5cm of cond] (end) {Penyatuan Alur Eksekusi};
  \draw[arrow] (start) -- (cond);
  \draw[arrow] (cond) -| node[near start, left, font=\bfseries, text=blue] {Terpenuhi} (t);
  \draw[arrow] (cond) -| node[near start, right, font=\bfseries, text=red] {Tidak Terpenuhi} (f);
  \draw[arrow] (t) |- (end);
  \draw[arrow] (f) |- (end);
\end{tikzpicture}
\caption{Visualisasi keputusan bercabang berdasarkan evaluasi Boolean.}
\label{fig:alur-boolean}
\end{figure}

\subsection{Menyusun Kondisi Kompleks dengan Operator}

Misalkan sebuah sistem diminta memverifikasi kelayakan pelamar kerja berdasarkan sejumlah persyaratan logika. Dalam industri \textit{Google} dan \textit{Mobile Banking}, algoritma penyaringan memproses banyak kondisi secara bersamaan yang menuntut penggunaan operator logika \textbf{AND (Konjungsi), OR (Disjungsi), dan NOT (Negasi)} \cite{python_docs_boolean}.

\begin{contoh}[Simulasi Prosedur Pengamanan Sistem]
Kasus \textit{Cyber Security}: akses brankas transfer dana diberikan hanya kepada pengguna yang memenuhi persyaratan berikut: pengguna memiliki alamat IP domestik (\textbf{ATAU} berstatus staf VIP) \textbf{SERTA} usianya tidak termasuk kelompok balita bebas rekening.

\begin{lstlisting}[style=pythonstyle, caption={Operasi Filter Keamanan dengan Operator Boolean}]
# Skema Filter Keamanan Finansial:
staf_vip = False
ip_domestik = True
usia_akun_tahun = 3

# Menyusun ekspresi evaluasi Boolean:
# Status evaluasi -> (True OR False) AND (NOT False)
# Status evaluasi -> (True) AND (True)
izin_lolos = (ip_domestik or staf_vip) and not (usia_akun_tahun < 12)

if izin_lolos:
    print("Akses brankas diberikan. Validasi sah.")
else:
    print("Akses ditolak.")
\end{lstlisting}
Melalui penguasaan hukum ekuivalensi (misalnya Hukum De Morgan), perancang basis data dapat menyederhanakan ekspresi \texttt{if} yang kompleks menjadi bentuk yang lebih ringkas dan efisien.
\end{contoh}

\section{Fondasi Kecerdasan Buatan dan Basis Data Relasional}

Logika matematika bukan sekadar instrumen untuk menganalisis percabangan \texttt{IF--ELSE} di perangkat lunak. Pada level makroskopis, arsitektur penyimpanan data (seperti Oracle dan PostgreSQL) maupun sistem kecerdasan buatan dibangun berbasis Logika Predikat Orde Satu (\textit{First-Order Logic}) \cite{mit_6042,rosen2019}.

\subsection{Pemrosesan Data Presisi dengan Aljabar Relasional (SQL)}

Sejumlah besar baris \textit{record} tabungan nasabah di bank tidak ditangani secara manual. Data tersebut diproses melalui \textbf{Aljabar Relasional} yang pada hakikatnya merupakan turunan langsung dari operasi \textit{Konjungsi (AND)} dan \textit{Disjungsi (OR)} serta kuantifikasi predikat dari Logika Matematika pada abad ke-19.

\begin{contoh}[Eksplorasi Kueri Bahasa Basis Data SQL]
Sebuah \textit{kueri} dirancang untuk mengambil profil tertentu: ``Tampilkan seluruh mahasiswa Fakultas MIPA yang umurnya di bawah 30 tahun dan tidak memiliki jejak skorsing kriminal.''

Sintaks SQL (\textit{Structured Query Language}) murni mengadopsi struktur Proposisi Biner di atas:
\begin{lstlisting}[language=SQL, basicstyle=\ttfamily\small, frame=single, caption={Analisis Klausa \texttt{WHERE} dengan Operator Logika Boolean}]
SELECT nama_lengkap, nim 
FROM Tabel_Mahasiswa_Nasional
WHERE (fakultas = 'MIPA') 
  AND (usia_tahun < 30) 
  AND NOT (status_skorsing = 'Kriminal');
\end{lstlisting}
Struktur kurung serta operator \texttt{AND}, \texttt{OR}, dan \texttt{NOT} memastikan presisi basis data dalam menyaring kandidat \cite{rosen2019}.
\end{contoh}

\subsection{Landasan Penalaran pada *Rules Engine* Kecerdasan Buatan (AI)}

Memasuki ranah kecerdasan buatan, ditemukan cabang *Sistem Pakar (Expert Systems)* yang menggunakan *Knowledge-based System Reasoning Engine*. Pada dunia kedokteran, sistem pakar bekerja menurunkan deduksi \textit{Modus Ponens} dari kumpulan fakta yang tersedia untuk menghasilkan kesimpulan yang menjadi dasar pengambilan keputusan \cite{mit_6042,saylor_cs202}.

\begin{contoh}[Mesin Inferensi Medis---Proses Inferensi Berbasis Aturan]

\textbf{Basis aturan (\textit{rule base}) dari protokol WHO:}\\
Aturan 1: \textbf{JIKA} $Pasien_{x}$ memiliki suhu $> 39$ Celsius \textbf{DAN} bercak kulit \textbf{MAKA} diagnosis Cacar Akut.\\
Aturan 2: \textbf{JIKA} $Pasien_{x}$ didiagnosis menderita Cacar Akut \textbf{DAN} usia $< 5$ tahun \textbf{MAKA} prioritas Karantina ICU.

\textbf{Fakta premis tambahan:} \\
Anak A memiliki suhu 40 Celsius, bercak kulit, dan usia 4 tahun.

\textbf{Proses inferensi (\textit{Reasoning Engine}):} 
Berdasarkan \textit{Modus Ponens}, mesin melakukan deduksi bertahap dan menghasilkan rekomendasi: \textbf{Anak A diprioritaskan untuk Karantina ICU rawat inap.} \\
Kemampuan menghubungkan fakta yang tersirat menjadi rangkaian sebab-akibat merupakan pencapaian \textit{Inference Logic} yang dirintis oleh Aristoteles 20 abad lalu \cite{stanford_cs103}.
\end{contoh}

\section{Gerbang Sirkuit Digital: Logika pada Keping Silikon}

Pertanyaan mendasar adalah bagaimana Aljabar Boolean yang dipahami melalui konsep dalam buku Boole dapat diimplementasikan dalam perangkat keras komputer.

Jawabannya berkaitan dengan perkembangan penting arsitektur fisika pada abad ke-20, yaitu penemuan \textbf{transistor} semikonduktor yang menerjemahkan nilai abstrak Boolean 1 (True) dan 0 (False) menjadi tegangan listrik, yaitu tegangan tinggi (misalnya 5 Volt) dan tegangan rendah (0 Volt) \cite{grinnell_logic,hackatronic}.

\subsection{Dasar Gerbang Logika (\textit{Logic Gates})}

Transistor berukuran nano disusun membentuk \textit{Logic Gates} yang mengimplementasikan operasi proposisional dasar AND, OR, dan NOT:

\begin{itemize}
    \item \textbf{Gerbang AND (Konjungsi):} Dua input (A, B). Output bernilai 1 hanya apabila kedua input bernilai 1. Jika salah satu input bernilai 0, output bernilai 0.
    \item \textbf{Gerbang OR (Disjungsi):} Dua input (A, B). Output bernilai 1 apabila minimal satu input bernilai 1.
    \item \textbf{Gerbang NOT (Inverter):} Satu input. Output membalik nilai input: input 1 menghasilkan output 0, input 0 menghasilkan output 1.
\end{itemize}

\begin{figure}[H]
\centering
\includegraphics[width=0.95\textwidth]{logic-gates}
\caption{Simbol standar gerbang logika: NOT (kiri), AND (tengah), OR (kanan) dalam arsitektur CPU \cite{hackatronic}.}
\label{fig:logic-gates}
\end{figure}

\subsection{Optimasi Rangkaian: Efisiensi Biaya melalui Aljabar Boolean}

Mengapa perancang perangkat keras perlu menguasai Aljabar Boolean dalam perancangan rangkaian digital?

Diberikan spesifikasi sirkuit: \textbf{$Y = (A \land B) \lor (A \land \neg B)$}. Tanpa analisis, rancangan awal memerlukan 2 Gerbang AND, 1 Gerbang NOT, dan 1 Gerbang OR---total 4 komponen \cite{levin_discrete}.

Insinyur yang menguasai Aljabar Boolean melakukan penyederhanaan:
$$ Y = A \land (B \lor \neg B) \equiv A \land \text{True} \equiv \textbf{A} $$

Ekspresi dapat disederhanakan menjadi A saja. Rangkaian hanya memerlukan satu kabel penghubung dari input A ke output Y, tanpa gerbang AND, OR, maupun NOT. Biaya material dapat dikurangi secara signifikan.
Ini merupakan penerapan Aljabar Proposisional: pengurangan \textbf{BOM (Bill of Materials)} melalui penyederhanaan ekspresi Boolean \cite{grinnell_logic,hackatronic}.

\section{Peta Jalur Materi: Taksonomi Buku Ajar}

Untuk mencapai kompetensi sarjana ilmu komputer, mahasiswa tidak hanya mempelajari tabel kebenaran dasar \cite{rosen2019,levin_discrete}.

Silabus buku ajar ``Logika Informatika'' ini disusun secara bertahap dengan pola spiral (Gambar~\ref{fig:peta-konsep}). Berikut adalah peta jalan materi yang akan dipelajari:

\begin{figure}[H]
\centering
\begin{tikzpicture}[
  node distance=0.6cm and 1cm,
  box/.style={rectangle, draw=black, rounded corners, fill=blue!10, text width=2.2cm, minimum height=0.7cm, align=center, font=\small},
  root/.style={rectangle, draw=black, rounded corners, fill=green!20, text width=3.2cm, minimum height=0.8cm, align=center, font=\small\bfseries},
  arrow/.style={->, thick, >=stealth}
]
  \node[root] (root) {Fondasi Logika Matematika};
  \node[box, below left=1.2cm and 0.5cm of root] (b3) {Bab 3 \& 4 \\Fondasi Proposisional Evaluatif};
  \node[box, below=1.2cm of root] (b4) {Bab 5 \& 6 \\ Aturan Inferensi Valid dan Sistem Pembuktian};
  \node[box, below right=1.2cm and 0.5cm of root] (b5) {Bab 7 \& 8\\Semesta Kuantifier Logika Predikat};
  \node[box, below=1.8cm of b4] (b6) {Bab 9 \& 10\\Induksi \& Rekursi Algoritma};
  \draw[arrow] (root) -- (b3);
  \draw[arrow] (root) -- (b4);
  \draw[arrow] (root) -- (b5);
  \draw[arrow] (b3) -- (b4);
  \draw[arrow] (b4) -- (b5);
  \draw[arrow] (b5) -- (b6);
  \draw[arrow] (b4) -- (b6);
\end{tikzpicture}
\caption{Spiral taksonomi hirarki logika: setiap tingkat mendukung pilar bab pada tingkat berikutnya.}
\label{fig:peta-konsep-bab2}
\end{figure}

\subsection{Pondasi Biner: Bab III--IV}
Dimulai dengan \textbf{Logika Proposisional} (nilai atomik) dan \textbf{Ekuivalensi/Tautologi} menggunakan tabel kebenaran serta hukum Aljabar Boolean (De Morgan). Kompetensi ini berguna untuk mendeteksi kesalahan logika dalam percabangan sintaks \texttt{IF--ELSE} perangkat lunak.

\subsection{Penarikan Konklusi: Bab V--VI}
Mahasiswa dilatih \textbf{Modus Ponens}, \textbf{Modus Tollens}, dan \textbf{pembuktian bertahap}. Materi mencakup Pembuktian Langsung (\textit{Direct Proof}), Pembuktian Kontradiksi, dan Kontraposisi.

\subsection{Logika Predikat: Bab VII--VIII}
Setelah menguasai proposisi atomik, mahasiswa mempelajari \textbf{Logika Predikat (kuantor $\forall$/$\exists$)}. Hal ini merupakan dasar perumusan kueri (\texttt{WHERE EXISTS}) basis data relasional dan pemodelan penalaran pada sistem pakar AI.

\subsection{Induksi dan Rekursi: Bab IX--X}
Bagian akhir mempelajari \textbf{Induksi Matematika} dan \textbf{Rekursi}. Mahasiswa diharapkan mampu membuktikan keabsahan algoritma iteratif (seperti \textit{sorting}/\textit{searching}) yang menggunakan perulangan.

Setiap bab merupakan prasyarat bagi bab berikutnya. Urutan materi hendaknya diikuti secara konsisten \cite{levin_discrete,rosen2019}.


% ============================================================
% AKTIVITAS PEMBELAJARAN
% ============================================================

\begin{aktivitas}
  \item \textbf{Analisis Kesalahan Inferensi:} Rancang sebuah argumen (2 premis + 1 kesimpulan) yang tidak valid secara struktur namun isi kalimatnya dapat bernilai benar secara faktual. Tukarkan argumen Anda dengan argumen invalid milik rekan kelas. Identifikasi mengapa aturan inferensinya cacat di mata logika formal.
  \item \textbf{Latihan Penyederhanaan Ekspresi Boolean:} Diberikan rumus sirkuit $(A \lor \text{False}) \land (B \land \text{True})$. Gunakan hukum identitas Aljabar Boolean untuk menyederhanakan ekspresi tersebut. Hitung pengurangan komponen gerbang dan estimasi efisiensi biaya material yang dapat dicapai.
  \item \textbf{Evaluasi Ekspresi Boolean dengan Python:} Jalankan skrip Python: \texttt{print((4 > 2) and not (3 == 3))}. Sebelum mengeksekusi, perkirakan apakah output bernilai \texttt{True} atau \texttt{False} berdasarkan evaluasi operan. Verifikasi dengan menjalankan kode.
\end{aktivitas}

% ============================================================
% LATIHAN DAN REFLEKSI
% ============================================================

\begin{latihan}
  \item Mempertimbangkan kontribusi George Boole yang secara fundamental mentransformasikan penalaran manusia menjadi operasi biner 1 (True) dan 0 (False), uraikan: Bagaimana dampaknya terhadap perkembangan Internet, Web modern, dan arsitektur RAM apabila konsep Aljabar Boolean tidak pernah dikembangkan?
  \item Jelaskan kaitan antara kesalahan inferensi (\textit{fallacious reasoning}) dengan kesalahan logika (\textit{logic bugs}) dalam penulisan kode perangkat lunak. Pada titik mana kesesatan ``Penegasan Konsekuen'' (\textit{affirming the consequent}) paling sering muncul dalam struktur percabangan \texttt{IF--ELSE}?
  \item Bandingkan: Apa perbedaan operasi yang dilakukan oleh kuantor predikat ($\forall$) dalam algoritma basis data SQL saat mengeksekusi klausa \texttt{WHERE} terhadap jutaan baris data dibandingkan dengan evaluasi \texttt{IF--ELSE} terhadap satu entitas?
  \item Uraikan \textbf{Fallacy of Affirming the Consequent} sebagai kesesatan inferensi yang tidak valid. Mengapa penalaran dari akibat ke sebab (B membuktikan A) dianggap tidak sah dalam logika formal?
  \item \textbf{Refleksi:} Di antara aplikasi utama logika matematika---pengendali alur sintaks perangkat lunak, filter query basis data/inferensi AI, atau rancangan sirkuit perangkat keras---bidang manakah yang paling sesuai dengan minat rekayasa Anda? Jelaskan alasannya.
\end{latihan}

% ============================================================
% ASESMEN
% ============================================================

\begin{asesmen}
\textbf{Instrumen Penilaian untuk Sub-CPMK Bab 2}

\textbf{A. Pemilihan Ganda Kritis}

\begin{enumerate}
  \item Kontribusi utama George Boole dalam \textit{The Laws of Thought} (1854) yang dapat dianggap sebagai fondasi bagi mikroprosesor transistor CPU di abad berikutnya adalah:
  \begin{enumerate}
    \item Merancang gerbang listrik 5 Volt.
    \item Merumuskan hukum inferensi Modus Ponens dalam bahasa awam.
    \item Mentransformasikan penalaran retorika menjadi operasi matematika biner (1/0) Boolean.
    \item Menemukan arsitektur chip motherboard.
  \end{enumerate}
  Jawaban: (c)

  \item Argumen: ``Jika Naga adalah mamalia, maka Naga betina bertelur di dasar samudera.'' Diberikan premis ``Naga adalah mamalia'' (bernilai benar dalam konteks fiksi). Status validitas argumen ini secara logika formal adalah:
  \begin{enumerate}
    \item Invalid, karena premis bersifat fiktif.
    \item Valid secara struktur, karena mematuhi aturan inferensi terlepas dari kebenaran faktual premis.
    \item Tidak dapat ditentukan karena variabel Naga belum didefinisikan dalam himpunan.
    \item Sesat pikir penegasan berbalik.
  \end{enumerate}
  Jawaban: (b)

  \item Diberikan percabangan \path{IF (Status_Dewa OR Armor_Emas) AND (NOT Terkena_Racun)}. Jika pemain terkena racun (\texttt{Terkena\_Racun = True}) tetapi memakai Armor Emas (\texttt{Armor\_Emas = True}), hasil evaluasi ekspresi Boolean untuk karakter tersebut adalah:
  \begin{enumerate}
    \item \texttt{True}, karena komponen OR terpenuhi.
    \item \texttt{False}, karena komponen \texttt{AND NOT Terkena\_Racun} bernilai \texttt{False}.
    \item \texttt{True}, karena Status Dewa menang.
    \item Tidak dapat ditentukan.
  \end{enumerate}
  Jawaban: (b)
\end{enumerate}

\textbf{B. Studi Kasus Evaluasi}

\begin{enumerate}
  \item Spesifikasi sirkuit menetapkan dua Gerbang AND yang disambungkan ke satu Gerbang NOT pada PCB. Gunakan hukum idempoten atau absorpsi Aljabar Boolean untuk menyederhanakan sirkuit tersebut. Laporkan pengurangan komponen dan estimasi efisiensi biaya pengadaan material (\textit{Bill of Materials}). Tunjukkan langkah penyederhanaan dan hasil akhir.
\end{enumerate}

\textbf{Rubrik penilaian}: Lihat Lampiran.
\end{asesmen}

% ============================================================
% CHECKLIST KOMPETENSI
% ============================================================

\begin{checklist}
  \item Saya mampu membedakan validitas inferensi aljabar (berdasarkan struktur) dengan kebenaran faktual premis dalam dunia empiris.
  \item Saya menguasai peran operator Boolean (AND, OR, NOT) sebagai dasar pengendali \textit{control flow} dalam bahasa pemrograman.
  \item Saya memahami bagaimana logika predikat orde pertama menjadi dasar algoritma query SQL relasional dan mesin inferensi sistem pakar AI.
  \item Saya mampu mencapai efisiensi rancangan gerbang logika (AND, OR) melalui penyederhanaan ekspresi Aljabar Boolean.
  \item Saya memahami alur silabus dari logika proposisional atomik hingga tautologi, ekuivalensi, dan rancangan sirkuit.
\end{checklist}

% ============================================================
% RANGKUMAN
% ============================================================

\begin{rangkuman}
Bab 2 menjelaskan landasan teoretis logika matematika dan penerapannya.

\textbf{Poin-poin utama:}
\begin{itemize}
  \item Logika bermula dari analisis bahasa awam (\textit{syllogistic}) Aristoteles, kemudian dikembangkan menjadi operasi matematika biner oleh \textbf{George Boole}.
  \item \textbf{Pembedaan validitas} memisahkan ketat antara kebenaran faktual dengan keabsahan struktur inferensi. Suatu argumen dapat valid secara formal meskipun premisnya bersifat fiktif, asalkan mematuhi aturan inferensi.
  \item Kompiler perangkat lunak modern menggunakan prinsip kalkulus Boolean untuk mengendalikan percabangan \texttt{IF--ELSE} dalam algoritma.
  \item Logika predikat (kuantor) menjadi dasar pencarian terstruktur pada basis data relasional (SQL) dan mesin inferensi sistem pakar AI.
  \item Pemetaan tegangan listrik High (1) dan Low (0) menerjemahkan hukum abstrak Boolean ke dalam gerbang logika sirkuit silikon (AND, OR, NOT). Penyederhanaan ekspresi Boolean dapat mengurangi biaya produksi material secara signifikan.
\end{itemize}

\textbf{Kata kunci}: Silogisme Aristotelian, Aljabar Boolean, anatomi argumen, validitas struktural vs kebenaran faktual, percabangan kode perangkat lunak, gerbang logika, aljabar relasional.
\end{rangkuman}

\ifSubfilesClassLoaded{
  \renewcommand{\bibname}{Daftar Pustaka}
  \bibliographystyle{plain}
  \bibliography{../references}
}{}
\end{document}
